% Input fragment for a manuscript. No existing manuscript is changed.
\section{A reproducible family of perimeter-control benchmark problems}
\label{sec:reusable-benchmark}

The benchmark separates regional regulation, queue-inclusive economic planning,
and deadline-constrained plan replacement. Each task supplies a common model,
initial information, command constraints, and evaluation interval. This
separation allows feedback regulators, predictive planners, and runtime
certification architectures to be compared on the properties they provide.

\subsection{Regional regulation}
Let $x=(n_{11},n_{12},n_{21},n_{22})$ denote the four destination inventories,
with $n_1=n_{11}+n_{12}$ and $n_2=n_{21}+n_{22}$. Define
$g_{ij}=n_{ij}G_i(n_i)/n_i$ on the positive-accumulation domain. The native
two-region model is
\begin{align}
 \dot n_{11}&=d_{11}-g_{11}+u_{21}g_{21},&
 \dot n_{12}&=d_{12}-u_{12}g_{12},\\
 \dot n_{21}&=d_{21}-u_{21}g_{21},&
 \dot n_{22}&=d_{22}-g_{22}+u_{12}g_{12}.
\end{align}
The task is to regulate both regional totals to 3000 vehicles, subject to
destination-inventory nonnegativity, regional storage, and common gate limits.
The performance measure is
\begin{equation}
 J_R=\int_0^{T_R}\left[(n_1(t)-3000)^2+(n_2(t)-3000)^2\right]\,dt,
 \qquad T_R=1.5\ \mathrm{h}.
\end{equation}
Settling is the first recorded time after which both regional errors remain
within 25 vehicles throughout the remaining experiment. The reproduced
almost-smooth CLF regulator follows Zhong et al.\ \cite{zhong2018robust}; the matched tracking MPC
uses the same plant, demand information, and gate limits. Time-varying-demand
tests supply the instantaneous demand to the CLF formula and the known demand
profile to the predictive planner. The benchmark reports these information
choices explicitly.

\subsection{Queue-inclusive economic planning}
The extended state is $x=(n_1,n_2,n_{11},n_{22},n_{10},n_{20},q)$, where $q$
is exterior waiting inventory and $n_{12}=n_1-n_{11}-n_{10}$,
$n_{21}=n_2-n_{22}-n_{20}$. With admission
$a=\min\{60q,12000u_{02}\}$ in vehicles per hour, the dynamics are
\begin{align}
 \dot n_1&=-g_{11}-u_{12}(g_{12}+g_{10})+u_{21}g_{21}
                    +d_{11}+d_{12}+d_{10},\\
 \dot n_2&=-g_{22}+u_{12}(g_{12}+g_{10})-u_{21}g_{21}
                    +a-u_{20}g_{20}+d_{22}+d_{21}+d_{20},\\
 \dot n_{11}&=-g_{11}+u_{21}g_{21}+d_{11},&
 \dot n_{22}&=-g_{22}+u_{12}g_{12}+\tfrac12a+d_{22},\\
 \dot n_{10}&=-u_{12}g_{10}+d_{10},&
 \dot n_{20}&=-u_{20}g_{20}+u_{12}g_{10}+d_{20},\\
 \dot q&=D(t)-a.
\end{align}
The cubic production, expressed in vehicles per hour, is
\begin{equation}
 G_i(n)=1.4877\times10^{-7}n^3-2.9815\times10^{-3}n^2+15.0912n.
\end{equation}
Physical rows require nonnegative destination inventories, $n_i\le10000$,
$q\le2000$, $0\le u_j\le1$, and
$|u_{21}-u_{12}|\le0.75$, $|u_{20}-u_{02}|\le0.75$.
Predictive perimeter control motivates the plan-optimization tool
\cite{geroliminis2013mpc}, while boundary-queue modeling motivates counting
exterior waiting \cite{haddad2017queues}. The economic objective counts both
moving and waiting vehicles:
\begin{equation}
 J_Q=\int_{t_a}^{T}(n_1+n_2+q)\,dt,
 \qquad t_a=\tfrac13\ \mathrm{h},\quad T=1\ \mathrm{h}.
\end{equation}
Every proposal has two constant command blocks: five minutes and the
remaining 35 minutes. The archived minute-20 information state is supplied
as a portable benchmark fixture. It is generated by the manuscript's
synthetic demand construction; the cubic curve follows the literature.

\subsection{Deadline-constrained replacement}
At the same activation time, an operator may require $n_1(T)\le b$ while
prioritizing Region 2 through
\begin{equation}
 J_D=\int_{t_a}^{T}(n_1+p n_2+q)\,dt.
\end{equation}
The benchmark varies $b$, $p$, sensing half-width $\delta$, and arrival
residual $\overline w$. Initial states and offered arrivals satisfy
\begin{equation}
 x(t_a)\in[\widehat x-\delta\mathbf1,\widehat x+\delta\mathbf1],
 \qquad D(t)=D_{\rm forecast}(t)+w(t),\quad |w(t)|\le\overline w.
\end{equation}
The planning tools receive identical physical and service rows. A plan
certificate covers this initial-state and measurable-residual family through
the absolute time $T$. A rejected replacement leaves the verified
continuation available through that same original expiry.

Predictive repair tests $P_\theta=(1-\theta)P+\theta P_B$ at
$\theta\in\{0,0.05,\ldots,1\}$, where $P_B$ is the independently verified
continuation. The first passing grid point minimizes modification on this
finite ordered grid. Every preceding grid point is checked; the procedure
uses no global monotonicity assumption. A task-aware scenario MPC provides
an additional planning tool, and its chosen commands undergo the same
independent certification procedure.

Results distinguish solver availability, finite-scenario feasibility,
uncertainty-family certification, and realized constraint satisfaction.
Traffic costs are compared within matched feasible cases. The released
protocol, parameter grids, exact commands, uncertainty contracts, and raw
results allow subsequent methods to reuse the same tasks and report both
physical/service outcomes and queue-inclusive traffic performance.
